% !TeX root = ncaa-manuscript-en.tex
\section{Related Work}\label{sec:related-work}

We distinguish three questions: how a model uses a supplied behavioral rule,
how that rule is expressed and checked as a constraint, and how the constraint
is delivered as learning input. The following studies address different parts
of this path; constraint compliance and preferences among admissible actions
are separate assessment targets.

\subsection{Reasoning-Based Driving and Reasoning Faithfulness}

Recent driving models connect visual observations, language reasoning,
and action planning. Alpamayo-R1 combines Chain of Causation (CoC)
reasoning with trajectory prediction and uses supervised fine-tuning
and reinforcement learning to improve reasoning--action
alignment~\cite{r1}. DriveLM connects perception, prediction, and
planning through graph-structured question--answer pairs, while
DriveVLM links scene description and analysis to hierarchical
planning~\cite{r2,r3}. Drive-R1 addresses reasoning--planning
misalignment and history-based shortcut learning through reward
design~\cite{r5}. These studies develop learnable representations of
driving decisions.\par

Reasoning faithfulness and reasoning--action consistency raise
related but distinct questions. Turpin et al.~\cite{Turpin2023} show that
chain-of-thought explanations can omit factors influencing an answer,
while Lanham et al.~\cite{Lanham2023} investigate whether predictions depend on the
stated reasoning. In driving, Drive the
Thoughts checks whether predicted trajectories match commitments
expressed in VLA reasoning~\cite{r18}. These assessments concern reasoning
faithfulness or reasoning--action consistency; constraint compliance with externally
supplied behavioral constraints is a separate assessment target.\par

\subsection{Traffic-Rule Formalization and Safety Enforcement}

Formalized traffic rules provide a foundation for expressing
admissible driving behavior. Esterle et al.~\cite{r6} encode traffic rules in
linear temporal logic (LTL), and Maierhofer et al.~\cite{Maierhofer2020} formalize interstate
rules for compliance assessment. Rulebooks
represent rule priorities and the ordering of behavioral outcomes,
whereas Responsibility-Sensitive Safety (RSS) defines mathematical
conditions for safe behavior and responsibility~\cite{Censi2019,ShalevShwartz2017}.
\par

Checking the formalized rules themselves is also an established
concern. Wan et al.~\cite{Wan2024} combine equivalence verification and model
checking to assess the semantic consistency and correctness of digital
traffic rules. DriveReg instead retrieves applicable
regulations and incorporates them into language-model
reasoning, supplying rules at inference~\cite{r7}. One of our controls shares
this principle: adapters trained without rule
text receive the rule only at evaluation
(Table~\ref{tab:setup-related-controls}).\par

Safety enforcement acts at another interface. Shielding restricts or
corrects actions to satisfy temporal-logic specifications during
learning or execution~\cite{Alshiekh2018}. Context-conditioned
Hamilton--Jacobi reachability incorporates geometry, friction, and
disturbances into a safety representation used by a control
filter~\cite{r19}. GuardSynth acts on training data and leaves all candidates
available at inference, making it complementary to mechanisms that filter
proposed actions.\par

\subsection{Natural-Language Formalization and Controlled Natural Language}

Natural-language formalization connects readable requirements with
machine-interpretable specifications. NL2TL abstracts atomic
propositions and learns to translate requirements into temporal logic~\cite{r8}.
nl2spec exposes correspondences between language fragments and
subformulas so users can revise ambiguous or erroneous
translations~\cite{r9}. Their central output is a formal
specification. GuardSynth's language-generation stage reverses this mapping:
it generates CNL from checked specifications and evaluates correspondence by
parsing sentences back with a supported fixed grammar to recover their fields
(Section~\ref{sec:method-checks}).\par

Controlled natural language (CNL) restricts language to support a
defined interpretation. Attempto Controlled English (ACE) connects
restricted English syntax with logical representations~\cite{r10}, while Kuhn~\cite{r13}
classifies CNLs by their properties and purposes.\par

\subsection{Incorporating Constraints into Learning and Generation}

Constraints can influence models through several mechanisms. Semantic
loss converts logical knowledge into a training objective that
encourages outputs to satisfy constraints~\cite{Xu2018}. We compare a
categorical adaptation of this objective with supplying rules as input text
(Section~\ref{sec:setup-comparisons}; Eq.~\ref{eq:setup-semantic-loss}). Giunchiglia
et al.~\cite{Giunchiglia2022} survey a broader range of approaches for combining deep
learning with logical constraints. Constrained
decoding acts during generation: Grid Beam Search enforces the inclusion
of specified words or phrases, while guided generation restricts outputs
using regular expressions and grammars~\cite{Hokamp2017,Willard2023}.
These mechanisms address constraint satisfaction at different levels;
a training penalty, a lexical requirement, and a grammatical restriction
provide different forms of control.\par

Training with task instructions or rationale supervision can improve
inference-time task performance. Wei et al.~\cite{Wei2022} instruction-tune
FLAN to improve zero-shot performance on unseen task types, and
Hsieh et al.~\cite{Hsieh2023} use model-generated rationales as auxiliary
training targets for smaller models. Our study evaluates this principle with
case-specific behavioral constraints and numerical parameters. Rule-paired
evaluation holds scenes and candidates fixed while changing the supplied rule,
and specification checking assesses construction and correspondence separately
from the model's choices.\par

Natural-language principles offer another route. Constitutional AI
uses principles to guide self-critique, revision, and AI-feedback
training~\cite{Bai2022}. Deliberative Alignment trains models to recall
and reason over a reusable safety policy; category-specific specifications
guide training-data generation and filtering~\cite{Guan2024}. The policy
text is then removed from the prompts used for supervised fine-tuning.
GuardSynth retains case-specific constraints and their numerical parameters
in both training and evaluation inputs. Our P1 condition adds
natural-language guards directly, while P2 uses EBLC-derived CNL at
the same learning interface. Their comparison examines constraint wording and
construction paths within this shared interface. Our preceding work showed
that natural-language guards improve action selection~\cite{r4}; GuardSynth
adds a checked construction path for such constraints.\par

\subsection{Positioning of This Work}

GuardSynth combines these foundations at the interface where behavioral
constraints become learning data. It binds constraint fields to
observations, supplied rules, or modeling assumptions; represents
supported constraints as Evidence-Bound Lifecycle Contracts (EBLC);
checks satisfiability and candidate admissibility within restricted logical models; and
generates CNL appended to the original CoC. Its contribution is this
connected development path. We also test whether constraints supplied during
training improve rule-dependent selection relative to adapters fine-tuned
without rule text and supplied the same rules only at evaluation. Checks are limited to supported profiles
and supplied evidence. Model evaluation distinguishes constraint violations
from selection of the preferred admissible candidate under a supplied
objective. This path complements learning objectives and runtime filters;
its evaluation uses supervised candidate selection rather than a comparison
of supervised and reinforcement learning.\par

Table~\ref{tab:related-work-scope} summarizes the main interfaces and
assessment targets of the selected research.\par

\begin{table}[htbp]
\caption{Main interfaces and assessment targets of related research}\label{tab:related-work-scope}
\footnotesize
\begin{tabular*}{\textwidth}{@{\extracolsep\fill}>{\raggedright\arraybackslash}p{0.26\textwidth}>{\raggedright\arraybackslash}p{0.29\textwidth}>{\raggedright\arraybackslash}p{0.36\textwidth}@{}}
\toprule
Research direction & Main interface & Main assessment or control target\\
\midrule
Reasoning-based driving~\cite{r1,r2,r3,r5} & Observations, reasoning, and planning & Driving performance and reasoning--action alignment\\\noalign{\vskip 3pt}
Reasoning assessment~\cite{Turpin2023,Lanham2023,r18} & Explanations, predictions, or trajectories & Reasoning faithfulness or reasoning--trajectory consistency\\\noalign{\vskip 3pt}
Traffic-rule formalization and verification~\cite{r6,Maierhofer2020,Censi2019,ShalevShwartz2017,Wan2024} & Informal rules and formal behavioral models & Rule semantics, priorities, correctness, or behavioral compliance\\\noalign{\vskip 3pt}
Regulation-augmented reasoning~\cite{r7} & Retrieved regulations and scene reasoning & Selection and use of applicable regulations\\\noalign{\vskip 3pt}
Shielding and safety filtering~\cite{Alshiekh2018,r19} & Proposed actions or control inputs & Satisfaction of specified safety conditions\\\noalign{\vskip 3pt}
Natural-language formalization~\cite{r8,r9} & Language requirements and temporal logic & Translation quality and, in nl2spec, interactive correction\\\noalign{\vskip 3pt}
Logical learning and constrained decoding~\cite{Xu2018,Giunchiglia2022,Hokamp2017,Willard2023} & Training objectives or generation procedures & Logical satisfaction, required lexical content, or output structure\\\noalign{\vskip 3pt}
Instruction tuning and rationale distillation~\cite{Wei2022,Hsieh2023} & Task instructions or rationale supervision & Instruction generalization or data-efficient distillation\\\noalign{\vskip 3pt}
Natural-language principle learning~\cite{Bai2022,Guan2024} & Principles, training examples, and policy reasoning & Alignment with supplied principles or safety specifications\\
\midrule
\textbf{GuardSynth} & \textbf{Source-bound contracts and CoC learning inputs} & \textbf{Supported contract properties, field--CNL correspondence, and separately measured constraint compliance and preferred selection}\\
\botrule
\end{tabular*}
\end{table}
